\documentclass[12pt]{scrartcl}

%Matthis Vorlage
\usepackage{../mathe_definition}
%wird in Matthis Volage auch gemacht aber das rafft mein texstudio nicht und highlighted Dinge falsch
\usepackage{amsmath}

%Ausgelagerter Code fürs schöne seitenlayout
\date{\today}
\newcommand{\Blattnummer}{%
	5%
}

\title{{\itshape Stochastische Prozesse}}
\module{Blatt 1}
\author{Philip Wunderlich \and Jonathan Langer}
%\matriculation{3737927\and 3750030}

\fancypagestyle{netter_header}{
	\renewcommand{\and}{\\}
	\fancyhead[L]{\MyAuthor}
	\fancyhead[C]{\MyTitle\\ \MyModule}
	\fancyhead[R]{\MyDate\\Seite \thepage{} von \pageref{LastPage}}
	\fancyfoot[C]{}
}
\pagestyle{netter_header}

%Abkürzungen
\newcommand{\einruecken}{\begin{addmargin}[25pt]{0pt}} 
	%LaTeX mag keine Umlaute :(
\newcommand{\ausruecken}{\end{addmargin}}
\newcommand{\R}{\ensuremath{\mathbb{R}}}
\newcommand{\Om}{\ensuremath{\Omega}}
\newcommand{\om}{\ensuremath{\omega}}
\newcommand{\A}{\ensuremath{\mathcal{A}}}
\newcommand{\B}{\ensuremath{\mathcal{B}}}
\newcommand{\C}{\ensuremath{\mathcal{C}}}
\newcommand{\WraumAnn}{\ensuremath{(\Om, \A, P)} ein Wahrscheinlichkeitsraum}
\newcommand{\MraumAnn}{\ensuremath{(E,\B)} ein Messraum}
\newcommand{\Xt}{\ensuremath{(X_t)_{t\in T}}}
\newcommand{\Yt}{\ensuremath{(Y_t)_{t\in T}}}
%\newcommand{\Xw}{\ensuremath{X_t(\omega)}}
%\newcommand{\Yw}{\ensuremath{Y_t(\omega)}}
%\newcommand{\Xq}{\ensuremath{X_q(\omega)}}
%\newcommand{\Yq}{\ensuremath{Y_q(\omega)}}
\newcommand{\E}{\ensuremath{\mathbb{E}}}
\newcommand{\1}{\ensuremath{\textbf{1}}}
\newcommand{\LwieLp}{\ensuremath{\mathcal{L}}}
% \L gibts anscheinend schon
\newcommand{\scriptL}{\ensuremath{\mathcal{L}}}
\newcommand{\scriptP}{\ensuremath{\mathcal{P}}}
\newcommand*\diff{\mathop{}\!\mathrm{d}}
\newcommand{\F}{\ensuremath{\mathcal{F}}}
\newcommand{\Dn}{\ensuremath{D_n}}
\newcommand{\cn}{\ensuremath{c_n}}
\newcommand{\ep}{\ensuremath{{e^{tc_n}}}}
\newcommand{\en}{\ensuremath{{e^{-tc_n}}}}
\newcommand{\N}{\ensuremath{\mathbb{N}}}
%\newcommand{\Xn}{\ensuremath{(X_n)_{n\in \N}}}
%\newcommand{\Yn}{\ensuremath{(Y_n)_{n\in \N}}}
\newcommand{\var}{\text{var}}
\newcommand{\textP}{\text{P}}
\newcommand{\K}{\ensuremath{\mathcal{K}}}

\begin{document}

\subsection*{Exercise 16}

Let $(\Om,\A,\textP)$ be some probability space, $\F$ be a filtration on it and $X$ be an $\R$-valued and $\F$-adapted process. Finally, let $\tau$ be an $\F$-stopping time. Show that $Y$ defined by
$$Y(\om):=
\begin{cases}
X_{\tau(\om)}(\om), & \text{if }\tau(\om)<\infty \\
0, & \text{else} 
\end{cases}$$
is $\F_{\tau}$-measurable.

\subsubsection*{Proof}
\einruecken

We recall from the lecture that $X_{\tau}$ is $\F_{\tau}$-measurable. Therefore, for $c\in \R$
$$X_{\tau}^{-1}((-\infty,c])\in \F_{\tau},$$
so for $c<0$ we get
$$Y^{-1}((-\infty,c])=X_{\tau}^{-1}((-\infty,c])\in \F_{\tau}.$$
Because $\{\tau=\infty\}\cap\{\tau\le n\}=\emptyset\in\F_n$ for all $n\ge 0$, we get by definition $\{\tau=\infty\}\in\F_{\tau}$, so
$$Y^{-1}((-\infty,c])=\{\tau=\infty\}\cup X_{\tau}^{-1}((-\infty,c])\in \F_{\tau}.$$
Together with $\B(\R)=\sigma((-\infty,c]:c\in \R)$, this now yields the $\F_{\tau}$-measurability of Y.

\ausruecken

\subsection*{Exercise 17}

Let $(\Om,\A,\textP)$ be some probability space, $\F$ be a filtration on it and $\sigma$ and $\tau$ be some $\F$-stopping times with $\sigma\le\tau$. For a fixed $A\in\F_{\sigma}$ show that
$$\sigma\1_A+\tau\1_{\Om\setminus A}$$
is also an $\F$-stopping time.

\subsubsection*{Proof}
\einruecken

Because $\sigma\le\tau$, we get for $n\ge 0$:
$$\{\tau\le n\}=\{\sigma\le n\}\cap\{\tau\le n\}.$$
For $A\in\F_{\sigma}$ we also have $\Om\setminus A\in\F_{\sigma}$, so together with
\begin{align*}
\{\sigma\1_A + \tau\1_{\Om\setminus A}\le n\}
&= (A\cap\{\sigma\le n\}) \cup (\Om\setminus A\cap\{\tau\le n\}) \\
&= (A\cap\{\sigma\le n\}) \cup (\Om\setminus A\cap\{\sigma\le n\})\cap\{\tau\le n\},
\end{align*}
we get $\{\sigma\1_A + \tau\1_{\Om\setminus A}\le n\}\in\F_n$ by definition of $\F_{\sigma}$ and the $\F$-stopping time property of $\tau$.
Hence,  $\sigma\1_A + \tau\1_{\Om\setminus A}$ is an $\F$-stopping time.

\ausruecken

\subsection*{Exercise 18}

Let $(Y_i)_{i\ge 1}$ be a sequence of i.i.d. random variables on some probability space $(\Om,\A,\textP)$ with $\textP(Y_1=1)=p$ and $\textP(Y_1=-1)=1-p$ for some $p\in (0,1)$ with $p\ne 1/2$.
Define
$$X_n:=\left(\frac{1-p}{p}\right)^{\sum_{i=1}^n Y_i},\qquad n\ge 0.$$
Find a suitable filtration $\F$ such that $X:=(X_n)_{n\ge 0}$ is an $\F$-martingale. Use this to investigate the ``probabilities of winning for players 1 and 2'' in the sense of Section 2.4.2., where the case $p=1/2$ was considered.

\subsubsection*{Solution}
\einruecken

Consider the natural filtration $\F = \F^{Y}$. For $n\ge 0$, we know that $\sum_{i=1}^n Y_i$ is $\F_n$-measurable from the lecture. Because
$$f:x\mapsto \left(\frac{1-p}{p}\right)^x=e^{\ln\left(\frac{1-p}{p}\right)\cdot x}$$
is continuous, we get the $\F_n$-measurability of $X_n = f\circ (\sum_{i=0}^n Y_i)$.
Again, because f is continuous, we know that $\left(\frac{1-p}{p}\right)^{Y_n}$ is independent of $\F_{n-1}=\sigma(Y_1,\dots,Y_{n-1})$, which yields
\begin{align*}
\E\left(\left(\frac{1-p}{p}\right)^{Y_n}\middle|\F_{n-1}\right)
&= \E\left(\frac{1-p}{p}\right)^{Y_n} \\
&= p\cdot\frac{1-p}{p}+(1-p)\cdot\frac{p}{1-p} \\
&= 1-p+p=1.
\end{align*}
Hence, by the $\F_n$-measurability of $X_{n-1}$, we get 
\begin{align*}
\E(X_n|\F_{n-1})
&= \E\left(X_{n-1}\cdot\left(\frac{1-p}{p}\right)^{Y_n}\middle|\F_{n-1}\right) \\
&= X_{n-1}\cdot\E\left(\left(\frac{1-p}{p}\right)^{Y_n}\middle|\F_{n-1}\right) \\
&= X_{n-1},
\end{align*}
so $X$ is an $\F$-martingale.

Now consider $A,B\in\N$. For
\begin{align*}
\tau_A &:= \inf\{n\ge 0: X_n=A\} \\
\tau_B &:= \inf\{n\ge 0: X_n=-B\}
\end{align*}
consider the combined stopping time $\tau_{A,B}:=\tau_A\land\tau_B$. We are interested in the probablity of $V_1:=\{\tau_{A,B}=\tau_A\}$.
Analogous to the proof in Section 2.4.2 of the lecture, we get
\begin{align*}
P(\{\tau_{A,B}=\infty\})
=\lim_{n\rightarrow \infty}P(\{\tau_{A,B}>n\cdot (A+B)\})
\le\lim_{n\rightarrow \infty}(1-p^{A+B})^n.
\end{align*}

With the notations $\tau=\tau_{A,B}, X_{\infty}=0$ and by the pointwise limit $\tau\land n \rightarrow \tau$, we have the almost sure convergence
$$X_{\tau\land n}\rightarrow X_{\tau}.$$
By the optional stopping theorem we also know
$$\E X_{\tau\land n}=\E X_0 = 1,$$
so the dominated convergence theorem together with $|X_{\tau\land n}|\le A+B$ yields $\E X_{\tau}=1$.
Considering
$$X_{\tau}(\om)=\begin{cases}
	a := \left(\frac{1-p}{p}\right)^A & \text{if } \tau_{A,B}(\om)=\tau_A(\om) \\
	b := \left(\frac{1-p}{p}\right)^{-B} & \text{if } \tau_{A,B}(\om)=\tau_B(\om) \\
	0 & \text{if } \tau_{A,B}(\om)=\infty
\end{cases}$$
and using $P(\tau_{A,B}<\infty)=1$ again, we now get
\begin{align*}
1 = \E X_\tau
&= a\cdot P(\{\tau_{A,B}=\tau_B\}) + b\cdot P(\{\tau_{A,B}=\tau_B\}) \\
&= a\cdot P(W_1)+ b\cdot (1-P(V_1)) \\
&= (a-b)\cdot P(V_1) + b,
\end{align*}
which finally yields
$$P(V_1)=\frac{1-b}{a-b}=\frac{1-\left(\frac{1-p}{p}\right)^{-B}}
{\left(\frac{1-p}{p}\right)^A - \left(\frac{1-p}{p}\right)^{-B}}.$$

\ausruecken
	
\end{document}